\chapter{Los números}
Siempre que estudiamos matemáticas, nos enseñan a usar los números, pero nunca nos enseñan qué son estos ``números''. Podríamos emplear una definición formal para ser rigurosos, sin embargo, no es la idea. Vamos despacio\dots

Los números, como siempre nos cuentan, aparecen más como una necesidad que como una consecuencia. Siempre se suele contar la historia de que para poder cuantificar cantidades de ``algo'' se crean estas entidades llamadas números. Estos podemos representarlos de diversas maneras, y expresarlos de muchas formas. Sin embargo, la primera intuición es utilizar los dedos. 

Si yo pregunto ¿Cuántos dedos tienes? Posiblemente, tu respuesta sea diez. No obstante ¿Qué es y qué representa la palabra diez? Responder esta pregunta puede ser extremadamente sencillo y trivial\footnote{Trivial: tan simple que se asume como cierto} o un completo lío. 

Los números son entidades que nos permiten contar, comparar, medir, modelar, entre muchas cosas más. Sin embargo, a día de hoy es casi un dogma nuestro sistema de numeración. 

\section{El sistema de numeración}

Imagine que quiere contar piedras en una bolsa, pero todavía no dispone de nombres ni de símbolos para expresar cuántas hay. Una forma sencilla de hacerlo sería asociar un dedo a cada piedra. Este recurso permite llevar un registro mientras la cantidad de piedras no supere el número de dedos disponibles. La idea importante es que, para contar y comunicar cantidades, necesitamos alguna forma de representarlas.

Un sistema de numeración es un conjunto de símbolos y reglas que permite representar números. Los números no deben confundirse con los símbolos que usamos para escribirlos: una misma cantidad puede tener representaciones diferentes según el sistema elegido. La elección de los símbolos y de la forma de agruparlos es una convención desarrollada históricamente, y no existe una única manera posible de hacerlo.

En la actualidad se utilizan distintos sistemas de numeración. Algunos de ellos son:
\begin{itemize}
  \item Sistema binario: utiliza dos dígitos, \(0\) y \(1\).
  \item Sistema octal: utiliza ocho dígitos, \(0,1,2,3,4,5,6,7\).
  \item Sistema decimal: utiliza diez dígitos, \(0,1,2,3,4,5,6,7,8,9\). Es el sistema de numeración de uso cotidiano más extendido.
  \item Sistema hexadecimal: utiliza dieciséis dígitos, \(0,1,2,3,4,5,6,7,\) \(8,9,A,B,C,D,E,F\), y es muy empleado en computación.
  \item Sistema sexagesimal: organiza las cantidades en grupos de sesenta. En la medición del tiempo todavía aparece esta forma de agrupamiento: por ejemplo, \(60\) segundos forman \(1\) minuto y \(60\) minutos forman \(1\) hora.
\end{itemize}

Existen muchos otros sistemas de numeración. Lo esencial es distinguir la cantidad que se quiere expresar de la representación elegida para escribirla. Al contar objetos, puede pensarse intuitivamente que se establece una correspondencia entre cada objeto y una marca, una palabra o un numeral; el sistema de numeración proporciona una manera compartida y organizada de realizar esa representación.

\subsection{El sistema decimal}

Este es el sistema que más va a utilizar, el sistema decimal. Este sistema consta de diez dígitos cada uno con un nombre:
\begin{itemize}
  \item \(0\): ``cero'' \(\rightarrow\) representa cantidad nula. Para el ejemplo de contar piedras significa que no hay piedras.
  \item \(1\): ``uno'' \(\rightarrow\) representa una unidad. Este número es importante, porque a partir de él y una operación llamada suma se obtiene el resto de cantidades.
  \item \(2\): ``dos'' \(\rightarrow\) representa dos unidades. Aunque esto pueda parecerle absurdo y trivial, matemáticamente hablando es una verdad. Para que usted (lector) no quede con un vacío existencial después de leer sobre números, puede pensarlo de la siguiente manera: si definimos un dedo (por ejemplo el índice) como unidad, luego podemos decir que tenemos \(2\) dedos índice. Uno en cada mano, porque cada uno representa una unidad, y dos unidades se representarán con el símbolo 2, o con dos dedos.
  \item \(3\): ``tres'' \(\rightarrow\) representa tres unidades. Análogo a lo que sucede con el número 2, el tres es el dos más una unidad. Y esto continúa así infinitamente.
\end{itemize}
Este sistema se llama decimal porque tiene en total diez dígitos individuales. Estos dígitos pueden combinarse para representar cantidades mayores a \(9\). Aquí, por ejemplo, el número \(17\) es una representación con dos números de una cantidad mayor a \(9\). Cuando llega a \(9\), en el sistema decimal, se le han agotado los símbolos disponibles, de esta manera puede intuir que esto es una gran limitación si pretende contar cualquier cosa. Para salvar esta problemática, se usa la combinación de varios dígitos, asignando la cantidad más significativa al número de la izquierda y la menor al número de la derecha.\footnote{Combinatoria se refiere a mezclar números donde su orden de posición es importante, es decir, 12 no es lo mismo que 21, ya que la cifra más significativa en el primer caso es 1 y en el segundo caso 2. Esto significa que el segundo número es mayor.} Utilizar más números agrupados para representar cantidades tan grandes como se necesite es la solución.

Los números a partir del los 90s tienen nombres muy bien definidos, pasando por las centenas (cien, doscientos, \dots), luego millares (mil, diez mil, cien mil, \dots) y así cada número tendrá un nombre asociado. 

De este modo, si se dice ciento cincuenta y ocho mil cuatrocientos treinta y tres, se sabe que se refiere específicamente al número 158433. Los nombres de los números pasando los millones, billones y trillones tienen nombres muy peculiares. Por ejemplo, el \href{https://es.wikipedia.org/wiki/G\%C3\%BAgol\#G.C3.BAgolplex}{gugol} (ver enlace\footnote{\href{https://es.wikipedia.org/wiki/G\%C3\%BAgol\#G.C3.BAgolplex}{\texttt{https://es.wikipedia.org/wiki/G\%C3\%BAgol\#G.C3.BAgolplex}}}). Estos números tan grandes no son muy usuales de nombrar, pero son divertidos de conocer.


\section{La suma (o adición)}

Ahora que se han introducido los números naturales, se puede estudiar una de las operaciones más elementales que realizamos con ellos: la adición, comúnmente llamada suma.

La adición toma dos números naturales y les asigna otro número natural. Esta idea puede expresarse simbólicamente como
\[
+\colon\mathbb{N}\times\mathbb{N}\to\mathbb{N}.
\]
No se preocupe si esta notación parece extraña a primera vista. Por ahora solo queremos expresar, de una manera compacta y precisa, que la adición recibe dos números naturales y produce un número natural. A continuación analizaremos cada parte de esta escritura.

\subsection{Los números naturales}

El símbolo \(\mathbb{N}\) representa el conjunto de los números naturales, es decir, los números que utilizamos principalmente para contar:
\[
1,2,3,4,5,\dots
\]
Dependiendo de la convención utilizada, el cero también puede considerarse un número natural. En este libro adoptaremos la convención
\[
\mathbb{N}=\{0,1,2,3,\dots\}.
\]
El conjunto de los números naturales es infinito: no existe un último número natural. Por grande que sea el número que elija, siempre habrá números naturales mayores que él.

\begin{note}
En matemáticas es importante prestar atención a palabras como \emph{siempre}, \emph{nunca}, \emph{todos} o \emph{ninguno}. Estas palabras expresan afirmaciones generales y, por lo tanto, deben cumplirse en todos los casos comprendidos por la afirmación.

Por ejemplo, decir que siempre existe un número natural mayor que otro significa precisamente que esto debe ser cierto sin importar qué número natural se haya elegido.
\end{note}

Volvamos ahora a la expresión
\[
+\colon\mathbb{N}\times\mathbb{N}\to\mathbb{N}.
\]

La parte
\[
\mathbb{N}\times\mathbb{N}
\]
indica que la operación recibe un par de números naturales. Por eso decimos que la adición es una \emph{operación binaria}: necesita dos números sobre los cuales actuar.

La parte
\[
\to\mathbb{N}
\]
indica que el valor obtenido también es un número natural. En otras palabras, al adicionar dos números naturales, obtenemos nuevamente un número natural.

Esta propiedad recibe el nombre de \emph{clausura}: los números naturales son cerrados respecto de la adición.

\subsection{Operación y operandos}

Supongamos que elegimos dos números naturales cualesquiera, \(a\) y \(b\). Podemos representarlos mediante el par
\[
(a,b).
\]
Al aplicarles la operación de adición, escribimos
\[
a+b.
\]
Los números \(a\) y \(b\) reciben el nombre de \emph{sumandos}. El símbolo \(+\) indica la operación de adición, y la expresión \(a+b\) representa la suma de ambos números.

Si además conocemos otro numeral \(c\) que representa ese mismo número, podemos escribir la igualdad
\[
a+b=c.
\]
Por ejemplo, si \(a=2\) y \(b=3\), entonces
\[
2+3=5.
\]
Aquí, \(2+3\) y \(5\) son dos expresiones que representan el mismo número natural. La primera muestra explícitamente la adición realizada; la segunda expresa su valor mediante un único numeral.

Veamos ahora cómo interpretar estas ideas en un caso concreto.

\section{Práctica de sumas con naturales}

Seguramente esté impresionado con la suma. Sin embargo, de poco le sirve si no sabe aplicarla. Para ello vamos a despejar un poco la mente de las definiciones rigurosas y vamos a sumar algunos números.

\begin{example}
Se le pide realizar la adición entre dos números naturales: \(302\) y \(22\). Como ambos pertenecen a \(\mathbb{N}\), puede aplicar la operación de adición:
\[
\begin{array}{r@{}}
    302 \\
    +22 \\
    \hline
    324 \\
  \end{array}
\]

Conviene detenerse un momento en qué significa matemáticamente esta cuenta. La adición es una operación que recibe dos números naturales y produce otro número natural. Simbólicamente,
\[
+\colon \mathbb{N}\times\mathbb{N}\to\mathbb{N}.
\]
Aquí, el símbolo \(\times\) indica el producto cartesiano\footnote{La multiplicación o producto entre números se estudiará más adelante.}; en este contexto, simplemente expresa que la operación recibe un par de números naturales.

En nuestro caso, ese par es \((302,22)\). Al aplicarle la operación de adición obtenemos
\[
302+22.
\]

Esta expresión ya representa la suma de \(302\) y \(22\). El paso siguiente consiste en \emph{evaluarla}, es decir, encontrar una representación más directa del número que denota:
\[
302+22=324.
\]

La igualdad anterior nos dice que \(302+22\) y \(324\) representan exactamente el mismo número natural. La diferencia está solamente en su forma de escritura: \(302+22\) muestra el número mediante una operación y sus operandos, mientras que \(324\) lo expresa directamente mediante un único numeral decimal.

Así, efectuar la cuenta no crea una nueva cantidad distinta de \(302+22\); permite reconocer y escribir de una manera más simple la cantidad que esa expresión ya representa.
\end{example}

Esta observación permite ver los números de una manera más flexible. Por ejemplo,
\[
302=210+92.
\]
Por lo tanto, en la suma anterior también podríamos sustituir \(302\) por esa expresión y escribir
\[
302+22=(210+92)+22.
\]

De este modo, un mismo número puede estar representado por muchas expresiones diferentes. El numeral \(302\), la expresión \(210+92\) y cualquier otra expresión que tenga el mismo valor representan una misma cantidad.

Esta es también una forma útil de interpretar las cuentas habituales: cuando se presenta una expresión como \(302+22\), la cantidad que representa ya está determinada. Lo que buscamos al efectuar la cuenta es evaluar esa expresión y encontrar una forma más directa de escribir su valor.

A continuación practicaremos algunas adiciones. Aunque se trata de una operación elemental y de uso frecuente, estos ejercicios servirán para familiarizarse con su notación y con la idea de evaluación que acabamos de introducir. Puede utilizar una calculadora para verificar los resultados.

\subsection*{Ejercicios}

En los siguientes ejercicios trabajaremos con distintas formas de representar un mismo número natural. Recuerde que expresiones diferentes pueden representar exactamente la misma cantidad.

\subsubsection*{Evaluar expresiones}

En cada caso, escriba el número natural representado por la expresión dada.

\begin{enumerate}
\item \(3+4\)

\item \(7+0\)

\item \(8+5\)

\item \(23+45\)

\item \(17+17\)

\item \(38+27\)

\item \(24+15+33\)

\item \(186+297\)

\item \(425+368\)

\item \(507+408\)

\item \(2347+1859\)

\item \(3999+2888\)

\item \(47328+96574\)
\end{enumerate}

En cada ejercicio, la expresión presentada ya representa un número natural. Su tarea consiste en encontrar una escritura equivalente mediante un único numeral.

\subsubsection*{Descomponer números}

Ahora realice el proceso contrario. Escriba cada número como la suma de dos números naturales.

\begin{enumerate}
\item \(7\)

\item \(12\)

\item \(25\)

\item \(40\)

\item \(83\)

\item \(125\)

\item \(302\)

\item \(1000\)
\end{enumerate}

No existe una única respuesta correcta. Por ejemplo,
\[
12=7+5,
\]
pero también
\[
12=10+2
\]
y
\[
12=6+6.
\]
Todas estas expresiones representan el mismo número.

Repita ahora algunos de los ejercicios anteriores, pero encuentre tres descomposiciones diferentes para cada número.

\subsubsection*{Completar expresiones equivalentes}

Complete los espacios de manera que ambos lados de la igualdad representen el mismo número.

\begin{enumerate}
\item \(8+4=\square\)

\item \(15=\square+6\)

\item \(20=13+\square\)

\item \(17+8=\square+10\)

\item \(30+12=40+\square\)

\item \(100=60+\square\)

\item \(125=100+\square\)

\item \(302=200+\square\)

\item \(302=210+\square\)

\item \(500+25=400+\square\)
\end{enumerate}

Observe que en algunos ejercicios no se pide simplemente evaluar una suma. Es necesario reconocer que una misma cantidad puede expresarse de distintas maneras.

\subsubsection*{Construir expresiones}

Encuentre una expresión que cumpla con cada condición.

\begin{enumerate}
\item Una suma de dos números naturales que represente \(10\).

\item Una suma de dos números naturales diferentes que represente \(20\).

\item Una suma de dos números naturales iguales que represente \(30\).

\item Tres números naturales cuya suma represente \(25\).

\item Dos expresiones diferentes que representen \(50\).

\item Tres expresiones diferentes que representen \(100\).

\item Una expresión que represente \(302\) sin utilizar el numeral \(302\).

\item Dos expresiones distintas que representen el mismo número natural, sin escribir directamente cuál es ese número.
\end{enumerate}

\subsection{Un algoritmo para evaluar sumas}

Hasta aquí hemos trabajado con sumas que pueden evaluarse mentalmente o mediante descomposiciones sencillas. Sin embargo, cuando los números crecen, resulta útil disponer de un procedimiento sistemático que permita encontrar el valor de una suma sin depender únicamente del cálculo mental.

El método que estudiaremos a continuación es el algoritmo usual de la adición en notación decimal. Este algoritmo no define qué es sumar: la adición ya ha sido presentada como una operación entre números. Su función es ayudarnos a evaluar, de manera ordenada y eficiente, expresiones como
\[
4387+2956.
\]
Para comprender por qué funciona, debemos recordar que la posición de cada cifra dentro de un numeral determina su valor. Por ejemplo,
\[
4387=4000+300+80+7
\]
y
\[
2956=2000+900+50+6.
\]
Por lo tanto,
\[
4387+2956
\]
puede escribirse como
\[
(4000+300+80+7)+(2000+900+50+6).
\]
Si agrupamos las cantidades que ocupan una misma posición decimal, obtenemos
\[
(4000+2000)+(300+900)+(80+50)+(7+6).
\]
Así, el problema puede resolverse trabajando por separado con unidades, decenas, centenas y unidades de millar. El algoritmo usual de la adición aprovecha precisamente esta estructura.

\subsubsection*{Alinear las posiciones}

Para aplicar el algoritmo, los numerales se escriben uno debajo del otro haciendo coincidir las cifras que ocupan la misma posición:
\[
\begin{array}{rrrrr}
 & 4 & 3 & 8 & 7\\
+& 2 & 9 & 5 & 6\\
\hline
\end{array}
\]
Las unidades quedan debajo de las unidades, las decenas debajo de las decenas, las centenas debajo de las centenas, y así sucesivamente.

Esta alineación es fundamental. No estamos colocando las cifras de este modo por una cuestión estética, sino porque queremos sumar entre sí las cantidades correspondientes a una misma posición decimal.

\subsubsection*{Comenzar por las unidades}

Observemos primero la columna de las unidades:
\[
7+6=13.
\]
Sin embargo, \(13\) no puede representarse mediante una única cifra en la posición de las unidades. Debemos descomponerlo:
\[
13=10+3.
\]
Es decir, \(13\) unidades equivalen a \(3\) unidades y \(1\) decena.

Por eso escribimos el \(3\) en la posición de las unidades y trasladamos la decena restante a la columna siguiente:
\[
\begin{array}{rrrrr}
 &   &   & 1 &  \\
 & 4 & 3 & 8 & 7\\
+& 2 & 9 & 5 & 6\\
\hline
 &   &   &   & 3
\end{array}
\]
A este procedimiento se lo conoce habitualmente como ``llevar uno''. Pero ese \(1\) no aparece de manera arbitraria: representa una decena que hemos obtenido al descomponer las \(13\) unidades.

\subsubsection*{Continuar con las decenas}

En la columna de las decenas tenemos ahora
\[
1+8+5=14.
\]
Como cada una de estas cantidades se encuentra en la posición de las decenas, esto representa en realidad
\[
10+80+50=140.
\]
Podemos descomponer las \(14\) decenas como
\[
14\text{ decenas}=1\text{ centena}+4\text{ decenas}.
\]
Por eso escribimos \(4\) en la posición de las decenas y trasladamos \(1\) centena a la columna siguiente:
\[
\begin{array}{rrrrr}
 &   & 1 & 1 &  \\
 & 4 & 3 & 8 & 7\\
+& 2 & 9 & 5 & 6\\
\hline
 &   &   & 4 & 3
\end{array}
\]

\subsubsection*{Continuar con las centenas}

Ahora evaluamos
\[
1+3+9=13
\]
en la columna de las centenas.
Estas \(13\) centenas equivalen a
\[
1\text{ unidad de millar}+3\text{ centenas}.
\]
Escribimos entonces \(3\) en la posición de las centenas y trasladamos \(1\) a la posición de los millares:
\[
\begin{array}{rrrrr}
 & 1 & 1 & 1 &  \\
 & 4 & 3 & 8 & 7\\
+& 2 & 9 & 5 & 6\\
\hline
 &   & 3 & 4 & 3
\end{array}
\]
Finalmente, en la columna de las unidades de millar obtenemos
\[
1+4+2=7.
\]
Así llegamos a
\[
\begin{array}{rrrrr}
 & 1 & 1 & 1 &  \\
 & 4 & 3 & 8 & 7\\
+& 2 & 9 & 5 & 6\\
\hline
 & 7 & 3 & 4 & 3
\end{array}
\]
y, por lo tanto,
\[
4387+2956=7343.
\]

\subsubsection*{¿Qué hace realmente el algoritmo?}

El algoritmo puede parecer una sucesión de reglas sobre cifras, pero en realidad no hace más que organizar descomposiciones como las que ya conocemos.

Cuando una columna produce una cantidad menor que \(10\), esta puede permanecer completamente en esa posición. Por ejemplo,
\[
3+4=7.
\]
Cuando produce una cantidad igual o mayor que \(10\), la descomponemos según el sistema decimal. Por ejemplo,
\[
8+7=15=10+5.
\]
Las \(5\) unidades permanecen en la columna de las unidades, mientras que las \(10\) unidades se reorganizan como \(1\) decena y pasan a la columna siguiente.

Lo mismo ocurre en cualquier posición:
\begin{align*}
  10&\text{ unidades}=1\text{ decena},\\
  10&\text{ decenas}=1\text{ centena},\\
  10&\text{ centenas}=1\text{ unidad de millar},
\end{align*}
y así sucesivamente.

Por esta razón, ``llevar'' una cantidad de una columna a otra no modifica el número representado. Únicamente cambia la forma en que distribuimos esa cantidad entre las distintas posiciones decimales.

\begin{note}
Es importante distinguir nuevamente entre la adición y el algoritmo utilizado para evaluarla.

La expresión
\[
4387+2956
\]
ya representa una suma. El procedimiento vertical que acabamos de estudiar es solamente una herramienta para encontrar una representación más directa de ese mismo número:
\[
4387+2956=7343.
\]

Existen otras formas de llegar al mismo valor. El algoritmo decimal es especialmente útil porque proporciona un procedimiento ordenado que puede aplicarse de la misma manera incluso cuando los números tienen muchas cifras.
\end{note}


\section{Los números enteros y la sustracción}

Hasta ahora hemos trabajado con números naturales y con la adición. Dentro de \(\mathbb{N}\), sumar dos números no presenta ningún inconveniente: la suma de dos naturales siempre vuelve a ser un número natural.

Sin embargo, aparece una dificultad cuando intentamos realizar ciertas operaciones que habitualmente llamamos restas. Por ejemplo,
\[
7-3=4
\]
no parece presentar ningún problema, pues \(4\in\mathbb{N}\). Pero si intercambiamos los números,
\[
3-7,
\]
ya no encontramos ningún número natural que pueda representar esa cantidad.

El conjunto \(\mathbb{N}\) resulta demasiado pequeño para expresar todas las cantidades que queremos considerar. Necesitamos ampliarlo.

\subsection{Los números opuestos}

Consideremos el número \(3\). Queremos introducir un número que, al sumarse con \(3\), produzca \(0\). A este número lo llamamos \emph{opuesto} de \(3\) y lo escribimos
\[
-3.
\]
Por definición de opuesto,
\[
3+(-3)=0.
\]
De manera general, el opuesto de un número \(a\) se representa mediante
\[
-a
\]
y cumple
\[
a+(-a)=0.
\]
El signo \(-\) que aparece en \(-a\) no representa aquí una resta entre dos números. Indica que estamos tomando el \emph{opuesto} de \(a\).

Por ejemplo,
\[
-5,\qquad -18,\qquad -302
\]
son números, de la misma manera que
\[
5,\qquad18,\qquad302
\]
también lo son.
Observe además que el opuesto de \(0\) es el propio \(0\):
\[
-0=0,
\]
pues
\[
0+0=0.
\]

\subsection{Los números enteros}

Los números negativos que acabamos de introducir no pertenecen al conjunto de los números naturales. Por ejemplo,
\[
-3\notin\mathbb{N}.
\]
donde el símbolo \(\notin\) significa que no pertenece.

Para trabajar conjuntamente con los números naturales, sus opuestos y el cero utilizamos el conjunto de los números enteros, representado mediante \(\mathbb{Z}\):
\[
\mathbb{Z}
=
\{\ldots,-3,-2,-1,0,1,2,3,\ldots\}.
\]
De este modo, los números naturales son un subconjunto de los enteros
\[
\mathbb{N}\subseteq\mathbb{Z}.
\]
Es decir, los números
\[
1,2,3,\ldots
\]
son los enteros positivos (y también naturales); los números
\[
-1,-2,-3,\ldots
\]
son los enteros negativos; y el \(0\) no es positivo ni negativo.

La adición puede ahora considerarse sobre este conjunto:
\[
+\colon\mathbb{Z}\times\mathbb{Z}\to\mathbb{Z}.
\]
Esto significa que podemos adicionar dos números enteros cualesquiera y obtener nuevamente un número entero.

Por ejemplo,
\[
5+(-2),\qquad
(-4)+7,\qquad
(-8)+(-3)
\]
son todas expresiones perfectamente válidas de adición entre números enteros.

\subsection{Dos usos del símbolo menos}

Llegados a este punto conviene hacer una distinción importante. El símbolo \(-\) aparece en matemáticas con dos funciones relacionadas, pero diferentes.

En la expresión
\[
-5,
\]
el símbolo \(-\) actúa sobre un único número: indica el opuesto de \(5\). En este caso podemos llamarlo \emph{signo menos} u \emph{operador de oposición}.

En cambio, en
\[
8-5,
\]
el símbolo aparece entre dos números e indica una sustracción.

Ambos usos están íntimamente relacionados. De hecho, la sustracción puede definirse utilizando únicamente la adición y el concepto de opuesto.

\subsection{La sustracción como suma del opuesto}

Dados dos números enteros \(a\) y \(b\), definimos
\[
a-b=a+(-b).
\]
Esta igualdad contiene la idea fundamental de la sustracción:
\[
\boxed{\text{Restar un número equivale a sumar su opuesto.}}
\]
Por ejemplo,
\[
8-5
\]
puede escribirse como
\[
8+(-5).
\]
Ambas expresiones representan exactamente el mismo número:
\[
8-5=8+(-5)=3.
\]
Del mismo modo,
\[
3-7=3+(-7)=-4.
\]
Ahora podemos comprender por qué la operación que antes no podía efectuarse dentro de los números naturales sí puede evaluarse dentro de los enteros:
\[
-4\in\mathbb{Z}.
\]
No hemos necesitado inventar un procedimiento completamente independiente de la adición. La sustracción queda definida mediante una operación que ya conocemos y el concepto de número opuesto.

\subsection{El signo acompaña al sumando}

Esta forma de escribir la sustracción resulta especialmente útil cuando aparecen varias cantidades.

Considere, por ejemplo,
\[
2-10.
\]
Utilizando la definición anterior podemos escribir
\[
2-10=2+(-10).
\]
Ahora vemos claramente cuáles son los dos sumandos:
\[
2
\qquad\text{y}\qquad
-10.
\]
El segundo sumando es el número \(-10\), no el número \(10\).

Como ahora tenemos una adición, podemos utilizar la propiedad conmutativa:
\[
2+(-10)=(-10)+2.
\]
Por lo tanto,
\[
2-10
=
2+(-10)
=
(-10)+2
=
-8.
\]
Aquí aparece una idea que será muy útil en cálculos posteriores: una vez que una expresión ha sido escrita como una suma de enteros, cada sumando conserva su propio signo.

Por ejemplo,
\[
7-4+2
\]
puede interpretarse como
\[
7+(-4)+2.
\]
Los sumandos son entonces
\[
7,\qquad -4,\qquad 2.
\]
Al trabajar de esta manera resulta más sencillo distinguir los números que participan de la operación y reorganizarlos cuando sea conveniente.

\begin{note}
No debe confundirse esto con afirmar que la sustracción es conmutativa.

Por ejemplo,
\[
7-4=3,
\]
mientras que
\[
4-7=-3.
\]
Lo que sí podemos reorganizar libremente son los sumandos una vez que la expresión ha sido escrita como una adición:
\[
7+(-4)=(-4)+7.
\]
\end{note}

\subsection{Restar un número negativo}

La definición
\[
a-b=a+(-b)
\]
también explica qué ocurre cuando el número que restamos ya es negativo.
Considere
\[
7-(-3).
\]
Restar \(-3\) significa sumar su opuesto:
\[
7-(-3)=7+\bigl(-(-3)\bigr).
\]
Pero el opuesto de \(-3\) es \(3\):
\[
-(-3)=3.
\]
Por lo tanto,
\[
7-(-3)=7+3=10.
\]
Esto explica la conocida regla según la cual ``dos signos menos se convierten en un más''. No se trata de una regla aislada que deba memorizarse: es una consecuencia de tomar el opuesto de un número que ya era negativo.

En general,
\[
-(-a)=a
\]
y, por lo tanto,
\[
a-(-b)=a+b.
\]

\subsection{Algunas consecuencias}

A partir de la adición y de los números opuestos podemos obtener inmediatamente varias relaciones útiles.

Para cualquier \(a\in\mathbb{Z}\),
\[
a+0=a,
\]
pues \(0\) es el elemento neutro de la adición.
Además,
\[
a+(-a)=0,
\]
pues \(-a\) es el opuesto de \(a\).

Como la sustracción se define mediante la suma del opuesto, tenemos
\begin{align*}
a-a=a+(-a)=0,\\
a-0=a+(-0)=a,
\end{align*}
y
\[
0-a=0+(-a)=-a.
\]
También,
\[
a-(-b)=a+b.
\]
Estas expresiones no necesitan aprenderse como reglas independientes. Todas ellas pueden deducirse a partir de una única idea:
\[
a-b=a+(-b).
\]



\section{Igualdades y ecuaciones}

Hasta ahora hemos utilizado expresiones como
\begin{gather*}
3+4=7 \quad
\text{o}\quad
8-5=3.
\end{gather*}
El símbolo \(=\) indica que las expresiones ubicadas a ambos lados representan el mismo número. Por esta razón decimos que forman una \emph{igualdad}.

En
\[
3+4=7,
\]
la expresión \(3+4\) se encuentra en el \emph{miembro izquierdo} de la igualdad y \(7\) en el \emph{miembro derecho}. Ambos miembros tienen el mismo valor.

Esta observación es importante: una igualdad no debe interpretarse simplemente como una indicación para realizar una cuenta. Por ejemplo,
\[
7=3+4
\]
es tan válida como
\[
3+4=7.
\]
También podemos escribir
\[
2+5=10-3,
\]
pues ambos miembros representan el número \(7\).

\subsection{Cuando aparece un valor desconocido}

Supongamos ahora que escribimos
\[
x+3=8.
\]
La letra \(x\) representa un número cuyo valor todavía desconocemos. Una igualdad que contiene uno o más valores desconocidos recibe el nombre de \emph{ecuación}.

Resolver la ecuación consiste en encontrar qué valor debe tomar \(x\) para que la igualdad sea verdadera.

En este caso buscamos un número que cumpla
\[
x+3=8.
\]
Podemos reconocer mentalmente que ese número es \(5\), pues
\[
5+3=8.
\]
Por lo tanto,
\[
x=5
\]
es una solución de la ecuación.

Sin embargo, no siempre será tan sencillo reconocer la solución a simple vista. Necesitamos una manera de transformar una ecuación sin alterar los valores que la hacen verdadera.

\subsection{Mantener una igualdad}

Considere una igualdad cualquiera
\[
a=b.
\]
Si ambos miembros representan el mismo número, al sumarles una misma cantidad continuarán representando el mismo número. Para cualquier entero \(c\),
\[
a+c=b+c.
\]
Por ejemplo, partamos de
\[
7=7.
\]
Si sumamos \(4\) a ambos miembros,
\[
7+4=7+4,
\]
obtenemos
\[
11=11.
\]
La igualdad se ha mantenido.

Lo mismo sucede si sumamos un número negativo. Si partimos nuevamente de
\[
7=7
\]
y sumamos \(-3\) a ambos miembros,
\[
7+(-3)=7+(-3),
\]
obtenemos
\[
4=4.
\]
Como restar \(3\) equivale a sumar \(-3\), esto suele escribirse simplemente como
\[
7-3=7-3.
\]

Por lo tanto, cuando se dice que ``restamos la misma cantidad en ambos miembros'', en realidad estamos utilizando una idea que ya conocemos: sumar el opuesto de esa cantidad en ambos lados.

\begin{note}
Para conservar una igualdad debemos realizar la misma transformación en ambos miembros.

No sería válido partir de
\[
7=7
\]
y escribir
\[
7+2=7+5,
\]
pues obtendríamos
\[
9=12,
\]
que es falso.
\end{note}

\subsection{Resolver una ecuación}

Volvamos ahora a
\[
x+3=8.
\]
Queremos conocer \(x\). Como el \(3\) está sumado a \(x\), podemos sumar su opuesto, \(-3\), en ambos miembros:
\[
x+3+(-3)=8+(-3).
\]
En el miembro izquierdo,
\[
3+(-3)=0,
\]
por lo que
\[
x+0=5.
\]
Como
\[
x+0=x,
\]
obtenemos finalmente
\[
x=5.
\]

No hemos ``movido el \(3\) al otro lado'' de la ecuación. Hemos realizado una misma operación sobre ambos miembros para mantener la igualdad.

Esta distinción será importante más adelante. Expresiones habituales como ``pasar un número al otro lado cambiando el signo'' son atajos útiles, pero ocultan la razón matemática por la que el procedimiento funciona.

\subsection{Otro ejemplo}

Consideremos
\[
x-4=9.
\]
Primero podemos interpretar la sustracción como una suma:
\[
x+(-4)=9.
\]
Queremos eliminar el sumando \(-4\). Para ello sumamos su opuesto, \(4\), en ambos miembros:
\[
x+(-4)+4=9+4.
\]
Como
\[
(-4)+4=0,
\]
queda
\[
x=13.
\]
Podemos comprobarlo sustituyendo \(x\) por \(13\) en la ecuación original:
\[
13-4=9.
\]

La igualdad es verdadera, por lo que \(13\) es efectivamente una solución.

\subsection{Una idea para recordar}

Resolver estas ecuaciones consiste, en gran medida, en utilizar números opuestos para aislar el valor desconocido.

Si tenemos
\[
x+a=b,
\]
podemos sumar \(-a\) en ambos miembros:
\[
x+a+(-a)=b+(-a),
\]
y obtenemos
\[
x=b+(-a).
\]
Es decir,
\[
x=b-a.
\]
La conocida idea de ``pasar \(a\) restando'' no es entonces una regla independiente: es una forma abreviada de sumar el opuesto de \(a\) en ambos miembros de una igualdad.


\section{Sumas de más de dos términos}

Hasta ahora hemos definido la adición como una operación binaria:
\[
+\colon\mathbb{Z}\times\mathbb{Z}\to\mathbb{Z}.
\]
Esto significa que la operación \(+\) recibe exactamente dos números cada vez que se aplica. Por ejemplo,
\[
3+5
\]
es una aplicación de la adición a los números \(3\) y \(5\).
Sin embargo, es habitual encontrar expresiones como
\[
3+5+7
\]
o incluso
\[
12+4+(-8)+15+(-2).
\]

A primera vista podría parecer que ahora estamos sumando tres, cuatro o cinco números simultáneamente. No es así. La adición continúa siendo una operación binaria: lo que hacemos es aplicarla varias veces.

\subsection{Sumar de a dos}

Considere la expresión
\[
3+5+7.
\]
Una manera de interpretarla es realizar primero
\[
3+5
\]
y después sumar \(7\) al número obtenido:
\[
(3+5)+7.
\]
Primero tenemos
\[
3+5=8,
\]
y luego
\[
8+7=15.
\]
Por lo tanto,
\[
(3+5)+7=15.
\]
También podríamos comenzar sumando \(5\) y \(7\):
\[
3+(5+7).
\]
En este caso,
\[
5+7=12
\]
y luego
\[
3+12=15.
\]
Obtenemos nuevamente el mismo número:
\[
3+(5+7)=15.
\]

Esto no ocurre por casualidad. La adición posee una propiedad llamada \emph{asociatividad}.

\subsection{La propiedad asociativa}

Para cualesquiera números enteros \(a\), \(b\) y \(c\),
\[
(a+b)+c=a+(b+c).
\]
Esta propiedad nos dice que, cuando únicamente estamos sumando, podemos cambiar la manera en que agrupamos los sumandos sin modificar el número representado.
Por ejemplo,
\[
(2+8)+5=2+(8+5).
\]
En ambos casos obtenemos
\[
15.
\]
Gracias a esta propiedad, normalmente omitimos los paréntesis y escribimos simplemente
\[
a+b+c.
\]
Por convención, esta escritura puede entenderse como aplicaciones sucesivas de la operación de adición. La propiedad asociativa garantiza que la forma de agruparlas no altera su valor.

Por esta razón podemos escribir sin ambigüedad expresiones como
\[
2+5+7+9.
\]

La operación sigue aplicándose de a dos números cada vez; simplemente hemos evitado escribir todos los paréntesis que serían necesarios para mostrar cada aplicación individual.

\subsection{Sumandos y términos}

En una expresión como
\[
2+5+7+9,
\]
los números
\[
2,\qquad5,\qquad7,\qquad9
\]
son los \emph{sumandos} de la adición.

También utilizaremos frecuentemente la palabra \emph{término} para referirnos a cada una de las partes principales que componen una expresión.

En una suma como la anterior, cada término es también un sumando. Más adelante encontraremos expresiones en las que un término puede ser algo más complejo que un único número.
Por ejemplo, en
\[
12+(3\cdot 4)+7,
\]
podemos considerar como términos principales
\[
12,\qquad 3\cdot4,\qquad7.
\]
Por ahora trabajaremos principalmente con términos que son números enteros.

\subsection{Expresiones con sustracciones}

La situación requiere un poco más de atención cuando aparecen signos de sustracción.
Considere
\[
123+3-2+(-20).
\]
Sabemos que toda sustracción puede escribirse como la suma del opuesto. Por lo tanto,
\[
123+3-2+(-20)
\]
puede reescribirse como
\[
123+3+(-2)+(-20).
\]
Ahora la expresión contiene únicamente adiciones.
Sus sumandos son
\[
123,\qquad3,\qquad-2,\qquad-20.
\]
Observe especialmente que uno de los sumandos es
\[
-2,
\]
no \(2\).

Una vez escrita toda la expresión como una suma, podemos aprovechar las propiedades de la adición para agrupar y ordenar los sumandos de la manera que resulte más conveniente.
Por ejemplo,
\[
123+3+(-2)+(-20)
\]
puede reorganizarse como
\[
123+(-20)+3+(-2).
\]
Luego podemos agrupar:
\[
\bigl(123+(-20)\bigr)+\bigl(3+(-2)\bigr),
\]
y evaluar:
\[
103+1=104.
\]
Por lo tanto,
\[
123+3-2+(-20)=104.
\]

\subsection{El signo forma parte del sumando}

Esta forma de interpretar una expresión permite aclarar algo que será muy importante durante los cálculos.
En
\[
7-4+9-12,
\]
podemos reemplazar las sustracciones por sumas de opuestos:
\[
7+(-4)+9+(-12).
\]
Ahora podemos identificar claramente los cuatro sumandos:
\[
7,\qquad-4,\qquad9,\qquad-12.
\]
Al reorganizar la suma, cada número conserva su signo:
\[
7+9+(-4)+(-12).
\]
Podemos incluso agrupar los positivos y los negativos:
\[
(7+9)+\bigl((-4)+(-12)\bigr).
\]
Así obtenemos
\[
16+(-16)=0.
\]
Esta manera de trabajar evita pensar que los signos se desplazan independientemente de los números. En una suma de enteros, números como
\[
-4
\qquad\text{y}\qquad
-12
\]
son sumandos completos.

\subsection{Una precaución con la sustracción}

La libertad para reagrupar aparece una vez que hemos expresado todo mediante adiciones.

La sustracción, por sí misma, no es asociativa.
Por ejemplo,
\[
(10-3)-2=7-2=5,
\]
mientras que
\[
10-(3-2)=10-1=9.
\]
Por lo tanto,
\[
(10-3)-2\neq10-(3-2).
\]
Cuando escribimos
\[
10-3-2,
\]
la convención habitual es interpretar las operaciones de izquierda a derecha:
\[
(10-3)-2.
\]
Sin embargo, podemos evitar esta dificultad reescribiendo primero las sustracciones:
\[
10-3-2
=
10+(-3)+(-2).
\]
Ahora sí tenemos únicamente una suma y podemos aplicar libremente las propiedades asociativa y conmutativa:
\[
10+(-3)+(-2)
=
10+\bigl((-3)+(-2)\bigr)
=
10+(-5)
=
5.
\]

\begin{note}
No debe confundirse la posibilidad de reorganizar una suma con la posibilidad de reorganizar arbitrariamente una expresión que contiene sustracciones.

Primero podemos transformar
\[
a-b+c-d
\]
en
\[
a+(-b)+c+(-d).
\]
Una vez hecho esto, tenemos una suma cuyos sumandos son
\[
a,\qquad -b,\qquad c,\qquad -d,
\]
y entonces sí podemos agruparlos u ordenarlos utilizando las propiedades de la adición.
\end{note}

\subsection{Una forma práctica de leer estas expresiones}

Ante una expresión como
\[
25-8+3-(-4)-10,
\]
puede resultar útil realizar primero una tarea puramente simbólica: identificar qué números se están sumando.

Reemplazamos cada sustracción por la suma del opuesto:
\[
25+(-8)+3+\bigl(-(-4)\bigr)+(-10).
\]
Como
\[
-(-4)=4,
\]
obtenemos
\[
25+(-8)+3+4+(-10).
\]
Ahora podemos leer la expresión como una suma de cinco términos:
\[
25,\qquad-8,\qquad3,\qquad4,\qquad-10.
\]
A partir de aquí podemos escoger una agrupación conveniente para evaluarla:
\[
25+(-10)+(-8)+3+4.
\]
Por ejemplo,
\[
\bigl(25+(-10)\bigr)+\bigl((-8)+3\bigr)+4
=
15+(-5)+4
=
14.
\]

El procedimiento puede resumirse de esta manera:
\begin{enumerate}
\item Reescriba cada sustracción como la suma de un opuesto.
\item Identifique los sumandos de la expresión.
\item Simplifique los dobles opuestos cuando aparezcan.
\item Agrupe u ordene los sumandos de una manera conveniente.
\item Evalúe las adiciones sucesivamente.
\end{enumerate}

No estamos introduciendo una nueva forma de sumar. Seguimos utilizando la misma operación binaria que ya conocemos, aplicada repetidas veces.


\section[Práctica con enteros y ecuaciones]{Ejercicios con números enteros y ecuaciones}

En los siguientes ejercicios utilizaremos conjuntamente las ideas desarrolladas hasta aquí. Recuerde que una sustracción puede escribirse como la suma del opuesto:
\[
a-b=a+(-b).
\]
También recuerde que, en una igualdad, cualquier transformación destinada a conservarla debe realizarse sobre ambos miembros.

\subsection*{Reescribir sustracciones}

Reescriba cada expresión utilizando únicamente adiciones. No es necesario evaluarlas todavía.
\begin{enumerate}
\item \(8-3\)

\item \(2-10\)

\item \(-5-7\)

\item \(14-(-5)\)

\item \(-8-(-3)\)

\item \(20-6-4\)

\item \(-15+8-(-3)\)

\item \(12-(-4)-9\)
\end{enumerate}
Por ejemplo,
\[
7-2+(-5)
\]
puede escribirse como
\[
7+(-2)+(-5).
\]
Una vez escrita de esta manera, resulta más sencillo identificar cada uno de los sumandos.

\subsection*{Evaluar expresiones con enteros}

Encuentre un único numeral que represente el valor de cada expresión.
\begin{enumerate}
\item \(12+(-7)\)

\item \((-9)+(-6)\)

\item \(8-15\)

\item \(14-(-5)\)

\item \(-8-(-3)\)

\item \(10-4+(-2)\)

\item \(-15+8-(-3)\)

\item \(45+(-28)-13+(-7)\)

\item \(-156+289-(-74)\)

\item \(-125+68-(-42)+(-89)-23\)
\end{enumerate}

Cuando resulte conveniente, reescriba primero todas las sustracciones como adiciones y reorganice luego los sumandos.

\subsection*{Completar igualdades}

Complete cada espacio con un número entero que haga verdadera la igualdad.

\begin{enumerate}
\item \(7+\square=12\)

\item \(15+\square=9\)

\item \(\square+8=3\)

\item \(10-\square=4\)

\item \(5-\square=12\)

\item \(-3+\square=8\)

\item \(12+(-5)=10+\square\)

\item \(20-8=\square+7\)

\item \(-4+9=12+\square\)

\item \(15-(-3)=20+\square\)
\end{enumerate}

Observe que aquí el objetivo no es simplemente realizar una operación: debe encontrar una cantidad que haga que ambos miembros representen el mismo número.

\subsection*{Resolver ecuaciones}

Encuentre el valor de \(x\) que hace verdadera cada ecuación. Escriba las transformaciones realizadas sobre ambos miembros.

\begin{enumerate}
\item

\[
x+4=11
\]

\item

\[
x+7=3
\]

\item

\[
x-5=12
\]

\item

\[
x-8=-2
\]

\item

\[
x+(-6)=10
\]

\item

\[
x-(-4)=9
\]

\item

\[
-7+x=5
\]

\item

\[
13+x=-4
\]

\item

\[
x+25=-18
\]

\item

\[
x-(-15)=-7
\]

\end{enumerate}

\subsection*{Comprobar soluciones}

Determine si el valor propuesto para \(x\) es solución de la ecuación. Para comprobarlo, sustituya \(x\) por el valor indicado y determine si obtiene una igualdad verdadera.

\begin{enumerate}
\item \(x+5=12\), con \(x=7\).

\item \(x-4=10\), con \(x=6\).

\item \(x+(-8)=-3\), con \(x=5\).

\item \(x-(-2)=9\), con \(x=7\).

\item \(-6+x=-10\), con \(x=-4\).

\item \(15+x=4\), con \(x=-11\).
\end{enumerate}


\section{Resumen del capítulo}

\begin{tcolorbox}[resumen,title=Sistemas de numeración]
Un sistema de numeración permite representar cantidades mediante símbolos y reglas.

En el sistema decimal utilizamos diez cifras:
\[
0,1,2,3,4,5,6,7,8,9.
\]
El valor de una cifra depende de la posición que ocupa. Por ejemplo,
\[
4387=4000+300+80+7.
\]
Por esta razón decimos que el sistema decimal es \emph{posicional}.
\end{tcolorbox}

\begin{tcolorbox}[resumen,title=Números naturales y enteros]
Los números naturales se representan mediante
\[
\mathbb{N}=\{0,1,2,3,\ldots\}.
\]
Los números enteros amplían este conjunto incorporando los números negativos:
\[
\mathbb{Z}
=
\{\ldots,-3,-2,-1,0,1,2,3,\ldots\}.
\]
Por lo tanto,
\[
\mathbb{N}\subseteq\mathbb{Z}.
\]
Todo número entero \(a\) posee un opuesto, representado por \(-a\), que cumple
\[
a+(-a)=0.
\]
\end{tcolorbox}

\begin{tcolorbox}[resumen,title=Adición]
La adición es una operación binaria:
\[
+\colon\mathbb{Z}\times\mathbb{Z}\to\mathbb{Z}.
\]
Dados dos números \(a\) y \(b\), su suma se representa mediante
\[
a+b.
\]
Los números que participan de la suma reciben el nombre de \emph{sumandos}.
Una expresión como
\[
302+22
\]
ya representa un número. Evaluarla consiste en encontrar otra representación equivalente de ese mismo número:
\[
302+22=324.
\]
Cuando aparecen más de dos sumandos, la adición continúa realizándose de a dos. La propiedad asociativa permite omitir los paréntesis cuando solo intervienen adiciones.
\end{tcolorbox}

\begin{tcolorbox}[resumen,title=Propiedades de la adición]
Para cualesquiera enteros \(a\), \(b\) y \(c\), la adición cumple:
  
\textit{Conmutatividad}
\[
a+b=b+a.
\]
\textit{Asociatividad}
\[
(a+b)+c=a+(b+c).
\]
\textit{Elemento neutro}
\[
a+0=a.
\]
\textit{Existencia del opuesto}
\[
a+(-a)=0.
\]

Estas propiedades permiten reorganizar y simplificar sumas sin modificar el número que representan.
\end{tcolorbox}

\begin{tcolorbox}[resumen,title=Sustracción]
La sustracción puede definirse mediante la adición y el opuesto:
\[
a-b=a+(-b).
\]
Por ejemplo,
\(
8-5=8+(-5).
\)

El símbolo \(-\) puede tener dos usos diferentes. En \(-5\) indica el opuesto de \(5\), mientras que en \(8-5\) indica una sustracción.

Una vez reescritas las sustracciones como sumas de opuestos, podemos trabajar utilizando las propiedades de la adición. Por ejemplo,
\[
7-4+2
=
7+(-4)+2.
\]
En esta última expresión los sumandos son \(7\), \(-4\) y \(2\).
\end{tcolorbox}

\begin{tcolorbox}[resumen,title=Igualdades y ecuaciones]
Una igualdad indica que las expresiones situadas a ambos lados de \(=\) representan el mismo número.

Por ejemplo,
\[
3+4=7
\]
y
\[
7=3+4
\]
expresan la misma igualdad.

Una ecuación es una igualdad que contiene uno o más valores desconocidos. Por ejemplo,
\[
x+3=8.
\]
Resolver una ecuación consiste en encontrar los valores que hacen verdadera la igualdad.

Para conservar una igualdad debemos realizar la misma transformación en ambos miembros. Por ejemplo,
\[
x+3=8
\]
puede transformarse sumando \(-3\) en ambos lados:
\[
x+3+(-3)=8+(-3),
\]
de donde
\[
x=5.
\]
La idea de ``pasar un término al otro lado cambiando el signo'' es una abreviatura de estas transformaciones y no una regla independiente.
\end{tcolorbox}

